\section{Baselines}
\label{sec:method-baseline}

Two baselines are used in this thesis. The two-stage OCR +
solver baseline
(Section~\ref{subsec:method-baseline-two-stage}) is the
non-joint alternative for visual Sudoku, against which the
end-to-end M\textsuperscript{3}N predictor is compared in
Chapter~\ref{chap:experiments}. The Forward-Backward marginal
classifier (Section~\ref{subsec:method-baseline-hmc-oracle}) is
the Bayes-optimal Hamming-loss decoder for symbolic HMC, used
as an oracle floor in the chain tightness experiments.

\subsection{Two-stage OCR + solver}
\label{subsec:method-baseline-two-stage}

The two-stage baseline is a deliberately simple alternative to
the joint M\textsuperscript{3}N predictor. It separates
perception from reasoning: a per-cell digit classifier first
reads the visual clues, and a hard-coded Sudoku solver then
completes the grid from those predicted clues. The baseline
serves two purposes in this thesis. First, it isolates the role
of the joint structured decision layer; any improvement of the
M\textsuperscript{3}N predictor over the two-stage pipeline is
attributable to end-to-end training. Second, the OCR component
uses the same backbone as the M\textsuperscript{3}N predictor's
unary feature extractor, so the comparison reduces to ``same
backbone, different decision layer''.

\paragraph{OCR digit classifier.}
A per-cell digit classifier maps each $28 \times 28$
MNIST-style image to a $K$-dimensional vector of class scores;
the predicted digit is the index of the maximum score. Only
clue cells of the visual Sudoku training split are fed to the
classifier --- blank cells are the all-zeros image
(Section~\ref{subsec:sudoku-visual}) and are dropped entirely.
With $N_{\mathrm{train}}$ training puzzles the classifier
therefore sees roughly $36\, N_{\mathrm{train}}$ clue images
per epoch. The classifier is trained with the cross-entropy
loss on 0-indexed digit labels, using the Adam optimizer with
the same hyperbolic learning-rate schedule as the structured
path (Section~\ref{subsec:method-training-optimizer}). All
hyperparameters live in the experiment YAML
(Section~\ref{sec:method-implementation}).

\paragraph{Hard-coded Sudoku solver.}
The solver is a standard depth-first backtracking search
augmented by naked-singles constraint propagation and the
Minimum Remaining Values (MRV) variable-ordering heuristic,
following Norvig~\cite{norvig2006sudoku} and the textbook
treatment of constraint-satisfaction problems
in~\cite{russell2020aima}. It either returns a complete
$9 \times 9$ grid satisfying all constraints, or signals
infeasibility when the predicted clues contradict each other.

\paragraph{Pipeline and metrics.}
At evaluation time the OCR classifier predicts a digit at
every clue cell; the predicted clues form a corrupted quiz,
which is passed to the solver; the solver's output is compared
cell-wise to the ground-truth solution. Three complementary
metrics are reported on every test split:
\begin{itemize}
    \item \emph{$0/1$ error.} The fraction of test puzzles
    whose final grid does not exactly match the ground-truth
    solution. Solver failures count as wrong. This is the
    headline metric for the two-stage system and is directly
    comparable to the M\textsuperscript{3}N predictor's $0/1$
    error on the same test split.
    \item \emph{Solver feasibility.} The fraction of puzzles
    where the solver returned some grid (even an incorrect
    one). Distinguishes ``OCR was confidently wrong''
    (feasible but incorrect) from ``OCR's clues are
    inconsistent with any Sudoku'' (infeasible).
    \item \emph{OCR clue error.} The fraction of clue cells
    where the OCR's argmax disagrees with the ground-truth
    clue digit. A pure component-level diagnostic of the OCR
    classifier, isolated from the solver.
\end{itemize}

\subsection{HMC oracle: Forward-Backward}
\label{subsec:method-baseline-hmc-oracle}

The HMC tightness experiments use the Forward-Backward marginal
classifier as the oracle floor. Given the known HMM parameters,
it computes per-cell posteriors $P(y_t \mid X)$ via the classic
Forward-Backward algorithm and predicts
$y_t = \argmax_y P(y_t = y \mid X)$, which minimizes expected
Hamming loss. No classifier with access to the same observations
can beat it; the structured-hinge and LP-M3N methods are
expected to approach but not surpass it. The Forward-Backward
decoder differs from Viterbi (used by structured-hinge for
loss-augmented inference at training time) in that it minimizes
Hamming rather than $0/1$ loss. The baseline is symbolic-only
by design; a visual oracle would require a learned per-cell
image classifier on top, which is not provided here.
